% Kapitel 3: Theoretische Fundierung
\section{Theoretische Fundierung}

\subsection{Human-Centred Design (ISO 9241-210)}

DIN EN ISO 9241-210~\cite{iso9241-210} beschreibt den nutzerzentrierten Gestaltungsprozess: verstehen, entwerfen, evaluieren, iterieren. Genau in dieser Reihenfolge sind wir vorgegangen -- von zehn Interviews über Personas und einen ersten Prototyp bis zu SUS-Tests und anschließenden UI-Anpassungen.

\subsection{Interaktionsprinzipien (ISO 9241-110)}

Die sieben Prinzipien nach DIN EN ISO 9241-110~\cite{iso9241} haben die Gestaltung geleitet, ohne dass für jede Kleinigkeit ein separates Kriterium gesucht wurde. Tabelle~\ref{tab:iso110} fasst die wichtigsten Bezüge zusammen; im weiteren Bericht werden einzelne Entscheidungen nur dort vertieft, wo sie für das Verständnis relevant sind.

\begin{table}[htbp]
\centering
\caption{Bezüge zu Interaktionsprinzipien nach ISO 9241-110}
\label{tab:iso110}
\begin{tabularx}{\textwidth}{@{}l X@{}}
\toprule
\textbf{Prinzip} & \textbf{Umsetzung in Ergonomia} \\
\midrule
Aufgabenangemessenheit & Kurze RAG-Antworten; nur relevanter Kontext \\
Selbstbeschreibungsfähigkeit & Begrüßung, Capability-Antwort, Vorschlagsfragen \\
Erwartungskonformität & Chat-Layout, bekannte UI-Muster (Composer, Sidebar) \\
Lernförderlichkeit & Vorschlagsfragen; Quellen zum Vertiefen \\
Steuerbarkeit & Neuer Chat; freier Themenwechsel \\
Individualisierbarkeit & Theme-Umschaltung (Blau/Orange); Session-Verwaltung \\
Fehlertoleranz & Fallback bei Wissenslücken; Disclaimer; Quellen zur Überprüfung \\
\bottomrule
\end{tabularx}
\end{table}

\subsection{Weitere Rahmenwerke}

Neben ISO nutzten wir mehrere etablierte Bezugsrahmen, je nachdem ob es um allgemeine Usability oder speziell um KI-Interaktion ging. Norman~\cite{norman2013} half vor allem bei Feedback und erkennbaren Bedienelementen -- etwa beim Streaming und den Quellenlinks. Nielsens Heuristiken~\cite{nielsen1994} waren praktisch beim Ladezustand und bei Fehlermeldungen. Für Chatbots mit LLM-Hintergrund griffen wir auf Amershi et al.~\cite{amershi2019} zurück; die Auswertung steht in Kapitel~\ref{sec:evaluation}. Der SUS~\cite{lewis2018} lieferte einen vergleichbaren Usability-Wert (Benchmark 68; Rohdaten im Anhang~\ref{app:sus}). Technisch basiert Ergonomia auf Retrieval-Augmented Generation~\cite{lewis2020rag}: Website-Inhalte werden abgerufen und dem Sprachmodell als Kontext mitgegeben, statt alles frei erfinden zu lassen.
